% Tests for latex-pro-code.sty, loaded on its own.
\documentclass{article}
\usepackage{latex-pro-code}
% regression: xcolor with options loaded after the package gives no clash
\usepackage[dvipsnames]{xcolor}
% regression: tcolorbox loaded after the package (it is only loaded at
% \begin{document} when missing)
\usepackage{tcolorbox}
\usepackage{lptest}
\makeatletter
\newcommand\pseudotmp{\jobname-pseudo.tmp}
\makeatother
\begin{document}
\lpdefined{python}{\python}
\lpdefined{py}{\py}
\lpdefined{pseudo}{\pseudo}
\lpdefined{endpseudo}{\endpseudo}
\python{code/sample.py}
\python*[firstline=1, lastline=2, caption={Part of a file}]{code/sample.py}
Inline \py{x -= f(x) / df(x)} and \py|a = {1: 2}|.
\begin{pythoncode}[caption=Typed in]
print("æøå")  # non-ASCII letters work with pdflatex
\end{pythoncode}

% \pseudo re-indents the file and turns operators into maths
\pseudo{code/indent.txt}
\lpsamefile{pseudo indents by keywords}{\pseudotmp}{code/indent-expected.txt}
\pseudoindentfalse
\pseudo*{code/indent.txt}
\lpsamefile{pseudoindentfalse keeps the indentation}{\pseudotmp}{code/indent-flat-expected.txt}
\pseudoindenttrue

% pseudocode typed in the .tex file goes through the same steps as \pseudo
\begin{pseudo}[caption=Typed in]
function search(A, x)
for i <- 1 to n do
if A_i = x then
return i
end if
end for
return 0
end function
\end{pseudo}
\lpsamefile{pseudo environment indents like the pseudo command}{\pseudotmp}{code/env-expected.txt}
% regression: the box was typeset while the verbatim catcodes were still
% active, so mathescape printed "\gets" instead of the arrow
\begingroup
\tcbset{latexpro-code/.append style={before upper=\xdef\lpcatbackslash{\the\catcode`\\}}}
\begin{pseudo}
x <- 1
\end{pseudo}
\endgroup
\lptrue{pseudo environment is typeset with normal catcodes}{\ifnum\lpcatbackslash=0 }
\begin{pseudo*}
if x then
y
end if
\end{pseudo*}
\lpsamefile{pseudo* environment}{\pseudotmp}{code/env-star-expected.txt}
% subscripts in pseudocode keep the code font: A_length, a_{n+1}
\begin{pseudo*}
x <- A_length + a_{n+1}
\end{pseudo*}
\lpsamefile{pseudo subscripts keep the code font}{\pseudotmp}{code/subscript-expected.txt}
\begin{pseudocode}
if x then
y
end if
\end{pseudocode}
\lpsamefile{pseudocode is the same environment}{\pseudotmp}{code/env-star-expected.txt}

% function names are set in small caps (\textsc) in a font that has them
\lpdefined{latexprofunctionfont}{\latexprofunctionfont}
\makeatletter
{\latexprocodefont\latexprofunctionfont
 \xdef\lpfnshape{\f@shape}%
 \xdef\lpfnfont{\expandafter\ifx\csname\f@encoding/\f@family/\f@series/\f@shape\endcsname\relax 0\else 1\fi}}
\makeatother
\lptrue{function names use the small-caps shape}{\ifnum\pdfstrcmp{\lpfnshape}{sc}=0 }
\lptrue{the code font has a small-caps shape}{\ifnum\lpfnfont=1 }
\makeatletter
\long\def\lpidentexpected{\iflatexpro@fn\global\latexpro@fnfalse\latexpro@functionstyle\fi}
\lptrue{name after def/function uses the function style}{\ifx\latexpro@ident\lpidentexpected}
\makeatother

% \variable{name}: a variable in running text, as in a pseudocode block
\lpdefined{variable}{\variable}
Text with \variable{lo}, \variable{A_i} and \variable{a_{n+1}^2}.
\lpsametext{variable is in the code font and text colour}
  {\variable{lo}}{{\latexprocodefont\color{latexpro@text}lo}}
% a subscript is set as maths, but letters and digits keep the code font
% (reported: A_length came out in maths italic)
\lpsametext{variable with a subscript keeps the code font}
  {\variable{A_i}}{$\latexprocodemath{A}_{\latexprocodemath{i}}$}
\lpsametext{variable with a word as subscript keeps the code font}
  {\variable{A_{length}}}{$\latexprocodemath{A}_{\latexprocodemath{length}}$}
\lptrue{variable subscript is not maths italic}
  {\sbox0{\variable{A_{length}}}\sbox2{$A_{length}$}\ifnum\numexpr\ifdim\wd0=\wd2 0\else 1\fi\relax=1 }
\lpsametext{variable with braces and superscript}
  {\variable{a_{n+1}^2}}{$\latexprocodemath{a}_{\latexprocodemath{n}+\latexprocodemath{1}}^{\latexprocodemath{2}}$}
\lptrue{variable is not in small caps (unlike method)}
  {\sbox0{\variable{merge}}\sbox2{\method{merge}}\ifdim\ht0<\ht2 }

% \method{name}: a function name in running text, styled as in the code
\lpdefined{method}{\method}
Text with \method{merge} in it.
\makeatletter
\lpsametext{method looks like a function name in code}
  {\method{merge}}{{\latexprocodefont\color{latexpro@function}\scshape merge}}
\makeatother
\lptrue{method uses small caps (taller than tt lowercase)}
  {\sbox0{\method{merge}}\sbox2{\ttfamily merge}\ifdim\ht0>\ht2 }

% // and /* */ comments in pseudocode. Regression: "/" was an operator
% (otherkeyword), which stopped listings from ever seeing "//" as a comment,
% and "*" in "/*" was turned into \cdot.
\pseudo{code/comments.txt}
\lpsamefile{pseudo comments: / coloured only outside comments, /* kept}{\pseudotmp}{code/comments-expected.txt}
\makeatletter
\newcommand\lpcomments[2]{% #1 listings style, #2 code: sets \lpsawcomment, \lpsawkw
  \begingroup
  \let\lporigcolor\color
  \def\lpkwname{latexpro@keyword}%
  \def\color##1{\def\lptmp{##1}\ifx\lptmp\lpkwname\gdef\lpsawkw{1}\fi\lporigcolor{##1}}%
  \gdef\lpsawkw{0}\gdef\lpsawcomment{0}%
  \sbox0{\lstinline[style=#1,commentstyle=\gdef\lpsawcomment{1}]{#2}}%
  \endgroup}
\makeatother
\lpcomments{latexpro-pseudo}{// a-b = c}
\lptrue{pseudo // starts a comment}{\ifnum\lpsawcomment=1 }
\lptrue{operators inside a pseudo comment are not coloured}{\ifnum\lpsawkw=0 }
\lpcomments{latexpro-pseudo}{x /* a-b */}
\lptrue{pseudo /* */ is a comment}{\ifnum\lpsawcomment=1 }
\lpcomments{latexpro-pseudo}{a-b}
\lptrue{operators in pseudo code are coloured}{\ifnum\lpsawkw=1 }
\lpcomments{latexpro-python}{x # a-b = c}
\lptrue{operators inside a Python comment are not coloured}{\ifnum\lpsawkw=0 }
\lpcomments{latexpro-python}{s = 1}
\lptrue{operators in Python code are coloured}{\ifnum\lpsawkw=1 }

% names with hyphens. Regression: in ADD-building-to-SKYLINE the "-" were
% operators, so only ADD and SKYLINE were function names (small caps) and
% the "-" were red. A "-" between two letters now joins one name; n-1 and
% a - b stay subtractions, and comments are left alone.
\pseudo{code/hyphen.txt}
\lpsamefile{pseudo keeps hyphenated names in one piece}{\pseudotmp}{code/hyphen-expected.txt}
\ExplSyntaxOn
\__latexpro_code_harvest_hyphen:n { code/hyphen.txt }
\lptrue{hyphenated~function~names~are~found~(lower~case)}
  { \str_if_eq:eeTF { \seq_use:Nn \l__latexpro_code_hyphen_names_seq { , } }
      { add-building-to-skyline } { \iftrue } { \iffalse } }
\bool_set_true:N \l__latexpro_code_hyphen_bool
\__latexpro_code_harvest:n { code/hyphen.txt }
\bool_set_false:N \l__latexpro_code_hyphen_bool
\lptrue{hyphenated~parts~are~not~function~names~on~their~own}
  { \str_if_eq:eeTF { \seq_use:Nn \l__latexpro_code_names_seq { , } }
      { MERGE } { \iftrue } { \iffalse } }
\ExplSyntaxOff
% #1 code (a line of a pseudocode block, after \pseudo's pre-processing):
% \lpsawfn is 1 when the function colour was used, \lpsawkw when the
% keyword (operator) colour was
\ExplSyntaxOn
\iow_new:N \g__test_colours_iow
\cs_new_protected:Npn \__test_colours_write:n #1
  {
    \iow_open:Nn \g__test_colours_iow { \c_sys_jobname_str -colours.tmp }
    \iow_now:Ne \g__test_colours_iow { \tl_to_str:n {#1} }
    \iow_close:N \g__test_colours_iow
  }
\NewDocumentCommand \lpcolours { +v }
  {
    \__test_colours_write:n {#1}
    \group_begin:
    \cs_set_eq:NN \lporigcolor \color
    \cs_set_protected:Npn \color ##1
      {
        \str_if_eq:nnT {##1} { latexpro@keyword } { \gdef\lpsawkw{1} }
        \str_if_eq:nnT {##1} { latexpro@function } { \gdef\lpsawfn{1} }
        \lporigcolor {##1}
      }
    \gdef\lpsawkw{0} \gdef\lpsawfn{0}
    \hbox_set:Nn \l_tmpa_box
      { \lstinputlisting [ style = latexpro-pseudo ] { \c_sys_jobname_str -colours.tmp } }
    \group_end:
  }
\ExplSyntaxOff
\lpcolours{function $\mbox{\latexprohyphenname{add-to-sky}}$(x)}
\lptrue{hyphenated name after function is a function name}{\ifnum\lpsawfn=1 }
\lpcolours{$\mbox{\latexprohyphenname{left-part}}$}
\lptrue{hyphenated variable is not a function name}{\ifnum\lpsawfn=0 }
\lptrue{hyphens in a variable name are not operators}{\ifnum\lpsawkw=0 }
\lpcolours{n-1}
\lptrue{minus before a digit is still an operator}{\ifnum\lpsawkw=1 }
\lpcolours{a - b}
\lptrue{minus between spaces is still an operator}{\ifnum\lpsawkw=1 }
\lpcolours{// function $\mbox{\latexprohyphenname{add-to-sky}}$}
\lptrue{hyphenated name in a comment is not a function name}{\ifnum\lpsawfn=0 }
\begingroup
\ExplSyntaxOn
\str_set:Nn \l_tmpa_str { add-to-sky }
\seq_clear:N \l__latexpro_code_hyphen_names_seq
\seq_put_right:NV \l__latexpro_code_hyphen_names_seq \l_tmpa_str
\ExplSyntaxOff
\lpcolours{x $\mbox{\latexprohyphenname{ADD-to-SKY}}$(1)}
\endgroup
\lptrue{a call of a hyphenated function is a function name}{\ifnum\lpsawfn=1 }
\lptrue{hyphens in a function name are not operators}{\ifnum\lpsawkw=0 }
\lpsametext{method with hyphens}
  {\method{ADD-building-to-SKYLINE}}{{\latexprocodefont\color{latexpro@function}\scshape ADD-building-to-SKYLINE}}
\lpsametext{variable with hyphens}
  {\variable{left-sum}}{{\latexprocodefont\color{latexpro@text}left-sum}}
\lpsametext{variable with hyphens and a subscript keeps one name}
  {\variable{left-sum_i}}{$\latexprocodemath{left-sum}_{\latexprocodemath{i}}$}
\lpsametext{variable subscript n-1 is still a subtraction}
  {\variable{a_{n-1}}}{$\latexprocodemath{a}_{\latexprocodemath{n}-\latexprocodemath{1}}$}

% Regression: a pseudo block with no \end{pseudo} ran on into the next
% block, and the "\begin{pseudo}[...]" line printed as code gave "Missing
% \endcsname". Now it is a clear error ("is an \end{pseudo} missing?") and
% the line is left out. The error is silenced here and counted instead.
\ExplSyntaxOn
\int_gzero:N \g_tmpa_int
\msg_redirect_name:nnn { latex-pro-code } { nested-begin } { none }
\cs_set_eq:NN \__test_orig_msg_error:nnVV \msg_error:nnVV
\cs_set_protected:Npn \msg_error:nnVV #1#2#3#4
  { \int_gincr:N \g_tmpa_int \__test_orig_msg_error:nnVV {#1} {#2} #3 #4 }
\ExplSyntaxOff
\pseudo{code/nested-begin.txt}
\lpsamefile{a begin pseudo line inside a block is left out}{\pseudotmp}{code/nested-begin-expected.txt}
\lptrue{a begin pseudo line inside a block is an error}{\ifnum\csname g_tmpa_int\endcsname=1 }
\ExplSyntaxOn
\cs_set_eq:NN \msg_error:nnVV \__test_orig_msg_error:nnVV
\msg_redirect_name:nnn { latex-pro-code } { nested-begin } { }
\ExplSyntaxOff

% pdflatex cannot load JetBrains Mono: the code font stays \ttfamily
\makeatletter
\lptrue{pdflatex keeps ttfamily as the code font}{\ifx\latexprocodefont\latexpro@defaultcodefont}
\makeatother
\lpdefined{wave height}{\latexprowaveheight}
\lpdefined{wave length}{\latexprowavelength}
% Regression: the wavy edges at a page break must fit together. The top
% of a wave on the bottom edge of one part (where it bulges down the most)
% is the bottom of a wave on the top edge of the next part, so the depth
% of the south cut plus the depth of the north cut is the same (4 wave
% heights) all along the edge. Before, both edges used the same wave and
% bulged out at the same places.
\makeatletter
\def\latexpro@waveh{1.8}\def\latexpro@wavel{14}
\newcommand\lpwavesum[1]{%
  \pgfmathparse{\latexpro@wavedepth{1}{#1}+\latexpro@wavedepth{-1}{#1}}%
  \pgfmathparse{abs(\pgfmathresult-7.2)<0.01}\ifnum\pgfmathresult=1 }
\newcommand\lpwavedepth[3]{%
  \pgfmathparse{abs(\latexpro@wavedepth{#1}{#2}-(#3))<0.01}\ifnum\pgfmathresult=1 }
\lptrue{wave edges fit together at 0pt}{\lpwavesum{0}}
\lptrue{wave edges fit together at 2pt}{\lpwavesum{2}}
\lptrue{wave edges fit together at 3.5pt}{\lpwavesum{3.5}}
\lptrue{wave edges fit together at 9pt}{\lpwavesum{9}}
\lptrue{wave edges fit together at 10.5pt}{\lpwavesum{10.5}}
\lptrue{wave edges fit together at 123.4pt}{\lpwavesum{123.4}}
% a quarter period in: the bottom edge dips (cut 3h up), the top edge
% reaches up (cut only h down), and three quarters in the other way round
\lptrue{bottom edge dips where the top edge reaches up}{\lpwavedepth{1}{3.5}{5.4}}
\lptrue{top edge reaches up where the bottom edge dips}{\lpwavedepth{-1}{3.5}{1.8}}
\lptrue{bottom edge bulges where the top edge dips}{\lpwavedepth{1}{10.5}{1.8}}
\lptrue{top edge dips where the bottom edge bulges}{\lpwavedepth{-1}{10.5}{5.4}}
% both waves start in the middle (depth 2h) at the west side, so parts of
% the same width line up
\lptrue{bottom wave starts at depth 2h}{\lpwavedepth{1}{0}{3.6}}
\lptrue{top wave starts at depth 2h}{\lpwavedepth{-1}{0}{3.6}}
\makeatother

% filename=false hides the file name in the header bar
\python[filename=false]{code/sample.py}
\pseudo[filename=false]{code/indent.txt}

% caption / title go to the header bar, the rest stays listings options
\ExplSyntaxOn
\__latexpro_code_opts:n { caption = {Newton's~method}, firstline=2, lastline=5 }
\lptrue{caption~is~taken~out~of~the~options}
  { \str_if_eq:eeTF { \tl_to_str:N \l__latexpro_code_caption_tl } { Newton's~method } { \iftrue } { \iffalse } }
\lptrue{other~options~go~to~listings}
  { \str_if_eq:eeTF { \clist_use:Nn \l__latexpro_code_lst_clist { ; } } { firstline=\c_left_brace_str 2\c_right_brace_str;lastline=\c_left_brace_str 5\c_right_brace_str } { \iftrue } { \iffalse } }
\cs_new_protected:Npn \__test_header:nn #1#2
  {
    \__latexpro_code_setup:nn {#1} {#2}
    \tl_gset_eq:NN \g_tmpa_tl \l__latexpro_code_caption_tl
  }
\__test_header:nn { } { file.txt }
\lptrue{header~shows~the~file~name~by~default}
  { \str_if_eq:eeTF { \tl_to_str:N \g_tmpa_tl } { file.txt } { \iftrue } { \iffalse } }
\__test_header:nn { filename = false } { file.txt }
\lptrue{filename=false~leaves~the~header~empty}
  { \tl_if_empty:NTF \g_tmpa_tl { \iftrue } { \iffalse } }
\__test_header:nn { filename = false , caption = C } { file.txt }
\lptrue{caption~still~shows~with~filename=false}
  { \str_if_eq:eeTF { \tl_to_str:N \g_tmpa_tl } { C } { \iftrue } { \iffalse } }
\__test_header:nn { filename } { file.txt }
\lptrue{bare~filename~key~shows~the~file~name}
  { \str_if_eq:eeTF { \tl_to_str:N \g_tmpa_tl } { file.txt } { \iftrue } { \iffalse } }
% [length=x]: the wave period of a split edge, x times the usual; default 1
\__latexpro_code_opts:n { length = 3 , firstline = 2 }
\lptrue{length~sets~the~wave~factor}
  { \str_if_eq:eeTF { \use:c { latexpro@wavefactor } } { 3 } { \iftrue } { \iffalse } }
\lptrue{length~is~not~passed~to~listings}
  { \str_if_eq:eeTF { \clist_use:Nn \l__latexpro_code_lst_clist { ; } }
      { firstline=\c_left_brace_str 2\c_right_brace_str } { \iftrue } { \iffalse } }
\__latexpro_code_opts:n { }
\lptrue{without~length~the~wave~factor~is~1}
  { \str_if_eq:eeTF { \use:c { latexpro@wavefactor } } { 1 } { \iftrue } { \iffalse } }
\__latexpro_code_opts:n { title = T }
\lptrue{title~is~the~same~as~caption}
  { \str_if_eq:eeTF { \l__latexpro_code_caption_tl } { T } { \iftrue } { \iffalse } }

% function names: name( and def/class/function NAME, keywords left out
\__latexpro_code_harvest:n { code/sample.py }
\lptrue{function~names~found~in~a~Python~file}
  { \str_if_eq:eeTF { \seq_use:Nn \l__latexpro_code_names_seq { , } }
      { newton,range,f,df,abs,RuntimeError } { \iftrue } { \iffalse } }
\cs_new_protected:Npn \__test_has:n #1
  {
    \str_set:Nn \l_tmpa_str {#1}
    \seq_if_in:NVTF \l__latexpro_code_names_seq \l_tmpa_str
      { \bool_set_true:N \l_tmpa_bool } { \bool_set_false:N \l_tmpa_bool }
  }
\__latexpro_code_harvest:n { code/indent.txt }
\__test_has:n { search }
\lptrue{function~names~found~in~pseudocode}
  { \bool_if:NTF \l_tmpa_bool { \iftrue } { \iffalse } }
\__latexpro_code_harvest:n { code/sample.py }
\__test_has:n { if }
\lptrue{keywords~before~a~parenthesis~are~not~function~names}
  { \bool_if:NTF \l_tmpa_bool { \iffalse } { \iftrue } }
\ExplSyntaxOff

% [span] is taken out of the listings options
\ExplSyntaxOn
\__latexpro_code_opts:n { span, firstline=2 }
\lptrue{span~sets~span=true}
  { \str_if_eq:VnTF \l__latexpro_code_span_tl { true } { \iftrue } { \iffalse } }
\lptrue{span~is~not~passed~to~listings}
  { \str_if_eq:eeTF { \clist_use:Nn \l__latexpro_code_lst_clist { ; } }
      { firstline=\c_left_brace_str 2\c_right_brace_str } { \iftrue } { \iffalse } }
\__latexpro_code_opts:n { firstline=2 }
\lptrue{without~span~it~is~auto}
  { \str_if_eq:VnTF \l__latexpro_code_span_tl { auto } { \iftrue } { \iffalse } }
\__latexpro_code_opts:n { span=false }
\lptrue{span=false~is~kept}
  { \str_if_eq:VnTF \l__latexpro_code_span_tl { false } { \iftrue } { \iffalse } }
\ExplSyntaxOff

% Long boxes. Regression: a box taller than a page was never split and ran
% off the bottom of the page. Now (span=auto) a box taller than a page
% breaks across pages by itself, a shorter one is kept in one piece.
\newcount\lpstart
\makeatletter
% \lpboxpage{label}: the page the label (set at the top of a box) is on
\def\lp@labelpage#1#2#3\@nil{#2}
\newcommand\lpboxpage[1]{\ifcsname r@#1\endcsname
  \expandafter\expandafter\expandafter\lp@labelpage\csname r@#1\endcsname\@nil
  \else 0\fi}
\newcommand\lplabelbox[1]{\tcbset{latexpro-code/.append style={before upper={\label{#1}}}}}
\makeatother
\clearpage \lpstart=\value{page}
\python{code/long.py}
\lptrue{a python box taller than a page breaks across pages}{\ifnum\value{page}>\lpstart}
\clearpage \lpstart=\value{page}
\pseudo{code/long.txt}
\lptrue{a pseudo box taller than a page breaks across pages}{\ifnum\value{page}>\lpstart}
\clearpage \lpstart=\value{page}
\begin{pythoncode}[caption=Long]
y_1 = 1 * 2
y_2 = 2 * 2
y_3 = 3 * 2
y_4 = 4 * 2
y_5 = 5 * 2
y_6 = 6 * 2
y_7 = 7 * 2
y_8 = 8 * 2
y_9 = 9 * 2
y_10 = 10 * 2
y_11 = 11 * 2
y_12 = 12 * 2
y_13 = 13 * 2
y_14 = 14 * 2
y_15 = 15 * 2
y_16 = 16 * 2
y_17 = 17 * 2
y_18 = 18 * 2
y_19 = 19 * 2
y_20 = 20 * 2
y_21 = 21 * 2
y_22 = 22 * 2
y_23 = 23 * 2
y_24 = 24 * 2
y_25 = 25 * 2
y_26 = 26 * 2
y_27 = 27 * 2
y_28 = 28 * 2
y_29 = 29 * 2
y_30 = 30 * 2
y_31 = 31 * 2
y_32 = 32 * 2
y_33 = 33 * 2
y_34 = 34 * 2
y_35 = 35 * 2
y_36 = 36 * 2
y_37 = 37 * 2
y_38 = 38 * 2
y_39 = 39 * 2
y_40 = 40 * 2
y_41 = 41 * 2
y_42 = 42 * 2
y_43 = 43 * 2
y_44 = 44 * 2
y_45 = 45 * 2
y_46 = 46 * 2
y_47 = 47 * 2
y_48 = 48 * 2
y_49 = 49 * 2
y_50 = 50 * 2
y_51 = 51 * 2
y_52 = 52 * 2
y_53 = 53 * 2
y_54 = 54 * 2
y_55 = 55 * 2
y_56 = 56 * 2
y_57 = 57 * 2
y_58 = 58 * 2
y_59 = 59 * 2
y_60 = 60 * 2
y_61 = 61 * 2
y_62 = 62 * 2
y_63 = 63 * 2
y_64 = 64 * 2
y_65 = 65 * 2
y_66 = 66 * 2
y_67 = 67 * 2
y_68 = 68 * 2
y_69 = 69 * 2
y_70 = 70 * 2
y_71 = 71 * 2
y_72 = 72 * 2
y_73 = 73 * 2
y_74 = 74 * 2
y_75 = 75 * 2
y_76 = 76 * 2
y_77 = 77 * 2
y_78 = 78 * 2
y_79 = 79 * 2
y_80 = 80 * 2
\end{pythoncode}
\lptrue{a long pythoncode environment breaks across pages}{\ifnum\value{page}>\lpstart}
\clearpage \lpstart=\value{page}
\begin{pseudo}[caption=Long]
x <- x + 1
x <- x + 2
x <- x + 3
x <- x + 4
x <- x + 5
x <- x + 6
x <- x + 7
x <- x + 8
x <- x + 9
x <- x + 10
x <- x + 11
x <- x + 12
x <- x + 13
x <- x + 14
x <- x + 15
x <- x + 16
x <- x + 17
x <- x + 18
x <- x + 19
x <- x + 20
x <- x + 21
x <- x + 22
x <- x + 23
x <- x + 24
x <- x + 25
x <- x + 26
x <- x + 27
x <- x + 28
x <- x + 29
x <- x + 30
x <- x + 31
x <- x + 32
x <- x + 33
x <- x + 34
x <- x + 35
x <- x + 36
x <- x + 37
x <- x + 38
x <- x + 39
x <- x + 40
x <- x + 41
x <- x + 42
x <- x + 43
x <- x + 44
x <- x + 45
x <- x + 46
x <- x + 47
x <- x + 48
x <- x + 49
x <- x + 50
x <- x + 51
x <- x + 52
x <- x + 53
x <- x + 54
x <- x + 55
x <- x + 56
x <- x + 57
x <- x + 58
x <- x + 59
x <- x + 60
x <- x + 61
x <- x + 62
x <- x + 63
x <- x + 64
x <- x + 65
x <- x + 66
x <- x + 67
x <- x + 68
x <- x + 69
x <- x + 70
\end{pseudo}
\lptrue{a long pseudo environment breaks across pages}{\ifnum\value{page}>\lpstart}
% a split box gets square corners and no border on the split edges: the
% jigsaw skin draws only the outer edges of each part
\clearpage
\makeatletter
{\tcbset{latexpro-code/.append style={before upper={\xdef\lpskin{\kvtcb@skin@current}}}}%
 \python{code/long.py}}
\def\lpwant{enhanced jigsaw}
\makeatother
\lptrue{code boxes use the jigsaw skin (open, square split edges)}{\ifx\lpskin\lpwant}
% span=false keeps even a long box in one piece
% (checked on the decision: where an over-full page is shipped out is not
% a stable thing to test)
\clearpage
\python[span=false]{code/long.py}
\ExplSyntaxOn
\lptrue{span=false~keeps~a~long~box~unbreakable}
  { \str_if_eq:VnTF \l__latexpro_code_span_tl { false } { \iftrue } { \iffalse } }
\ExplSyntaxOff
\clearpage
\python{code/long.py}
\ExplSyntaxOn
\lptrue{a~box~taller~than~a~page~is~made~breakable}
  { \str_if_eq:VnTF \l__latexpro_code_span_tl { true } { \iftrue } { \iffalse } }
\ExplSyntaxOff
\python{code/sample.py}
\ExplSyntaxOn
\lptrue{a~short~box~stays~unbreakable}
  { \str_if_eq:VnTF \l__latexpro_code_span_tl { false } { \iftrue } { \iffalse } }
\ExplSyntaxOff
% [span] still lets any box break
\clearpage \lpstart=\value{page}
\python[span]{code/long.py}
\lptrue{python span breaks across pages}{\ifnum\value{page}>\lpstart}
% a short box that does not fit at the bottom of a page moves to the next
% page whole (it is not split); with [span] it is split.  Regression: after
% one breakable box every later box was breakable too (\tcolorbox without a
% group kept the option), so these come after the [span] boxes above.
\clearpage \lpstart=\value{page}
\noindent\rule{1pt}{\dimexpr\textheight-10\baselineskip}\par
{\lplabelbox{lp-short}\python{code/sample.py}}
\lptrue{a short box moves to the next page whole}
  {\ifnum\lpboxpage{lp-short}=\numexpr\lpstart+1\relax}
\clearpage \lpstart=\value{page}
\noindent\rule{1pt}{\dimexpr\textheight-10\baselineskip}\par
{\lplabelbox{lp-short-span}\python[span]{code/sample.py}}
\lptrue{with span a short box is split at the page bottom}
  {\ifnum\lpboxpage{lp-short-span}=\lpstart}
% span=false in the pseudo environment, right after a [span] box: kept
% whole (reported: span=false did not work, because of the same leak)
\clearpage \lpstart=\value{page}
\noindent\rule{1pt}{\dimexpr\textheight-10\baselineskip}\par
{\lplabelbox{lp-pseudo-false}%
\begin{pseudo}[span=false]
function f(x)
x <- 1
x <- 2
x <- 3
x <- 4
x <- 5
x <- 6
x <- 7
return x
end function
\end{pseudo}}
\lptrue{pseudo environment with span=false moves to the next page whole}
  {\ifnum\lpboxpage{lp-pseudo-false}=\numexpr\lpstart+1\relax}
\end{document}
